\documentclass[11pt]{article}

\usepackage[margin=1in]{geometry}
\usepackage[T1]{fontenc}
\usepackage[utf8]{inputenc}
\usepackage{lmodern}
\usepackage{microtype}
\usepackage{amsmath,amssymb,amsthm,mathtools}
\usepackage{booktabs}
\usepackage{enumitem}
\usepackage{tabularx}
\usepackage[hidelinks]{hyperref}
\usepackage[nameinlink,capitalise]{cleveref}

\newtheorem{theorem}{Theorem}[section]
\newtheorem{proposition}[theorem]{Proposition}
\newtheorem{lemma}[theorem]{Lemma}
\newtheorem{corollary}[theorem]{Corollary}
\theoremstyle{definition}
\newtheorem{definition}[theorem]{Definition}
\newtheorem{assumption}[theorem]{Assumption}
\newtheorem{example}[theorem]{Example}
\theoremstyle{remark}
\newtheorem{remark}[theorem]{Remark}

\newcommand{\dd}{\mathrm d}
\newcommand{\cA}{\mathcal A}
\newcommand{\cC}{\mathcal C}
\newcommand{\cD}{\mathcal D}
\newcommand{\cE}{\mathcal E}
\newcommand{\cF}{\mathcal F}
\newcommand{\cG}{\mathcal G}
\newcommand{\cH}{\mathcal H}
\newcommand{\cM}{\mathcal M}
\newcommand{\cO}{\mathcal O}
\newcommand{\cR}{\mathcal R}
\newcommand{\cS}{\mathcal S}
\newcommand{\cU}{\mathcal U}
\newcommand{\cV}{\mathcal V}
\newcommand{\cW}{\mathcal W}
\newcommand{\MTT}{Modal Triplet Theory}
\newcommand{\Tr}{\operatorname{Tr}}
\newcommand{\tr}{\operatorname{tr}}
\newcommand{\Ker}{\operatorname{Ker}}
\newcommand{\BRST}{\mathrm{BRST}}
\newcommand{\SCFT}{\mathrm{SCFT}}

\title{\vspace{-1em}
Modal Triplet Theory and Perturbative String Theory:\\[0.25em]
\large A Conditional Encoding Contract and the \(q=79\) Worldsheet Boundary}
\author{Peter Nero}
\date{Version 2, July 2026}

\begin{document}
\maketitle

\begin{abstract}
This paper asks when Modal Triplet Theory (MTT) realizes a perturbative string
background.  A target metric, two-form, dilaton, and gauge connection are not
yet a string theory: one must also supply a two-dimensional quantum field
theory, its ghost and BRST systems, criticality, spin structures and GSO
projection, modular amplitudes, sewing and factorization, anomaly and global
\(B\)-field data, and an infrared prescription.  We package these ingredients
as a typed perturbative-string record and define a partial map from an upper
MTT state to that record.  The main theorem is conditional: if one selected
MTT state emits every required row, the lower record is string-consistent at
the declared order and genus, and the MTT observable map factors through the
standard worldsheet amplitude map, then MTT realizes that perturbative string
sector on the stated domain.  The theorem is an exact encoding statement, not
a derivation of string theory from saturation or projection alone.  We also
separate target-space beta-function equations from exact conformal
invariance, cohomological anomaly cancellation from the differential Bianchi
identity, central-charge arithmetic from model selection, and circle-spectrum
symmetry from full quantum T-duality.  The current selected \(q=79\) heterotic
program supplies five of twelve worldsheet rows and partially constructs two
more.  The physical non-pullback visible--hidden bundle, exact infrared
\((0,2)\) SCFT, analytic GSO and modular characters, quantum BV data, tadpole
control, and all-genus or nonperturbative completion remain open.
\end{abstract}

\section*{Revision note: Version 2}
\addcontentsline{toc}{section}{Revision note: Version 2}
\begin{description}[leftmargin=2.5cm,style=nextline]
\item[Supersedes]
Version 2 supersedes Version 1 and its claim that a bounded MTT projection had
already produced the Polyakov/RNS theory, conformal consistency,
Green--Schwarz anomaly cancellation, compactification, exact dualities, and
ultraviolet finiteness.

\item[Reason]
The former argument moved directly from ten-dimensional target fields to a
complete quantum worldsheet.  It conflated bosonic and heterotic
central-charge formulas, treated BRST nilpotence as synonymous with the
leading target beta functions, replaced a differential Bianchi identity by
equality of characteristic numbers, and assumed modular invariance in the
very step where it was needed.  It also imported withdrawn uniqueness and
stabilization claims from adjacent MTT papers.

\item[Resolution]
This revision defines the full lower string record and a typed MTT source map.
It proves the resulting conditional realization theorem, states a
fixed-order residual-transfer proposition, checks the standard heterotic
central-charge arithmetic, and records the current twelve-row \(q=79\)
worldsheet cutset.  T-duality, anomaly cancellation, fixed-genus ultraviolet
inheritance, and target-space reduction are restricted to the hypotheses
that actually support them.

\item[Retained content]
The ten-dimensional fields \((G,B,\Phi,A)\), sigma-model pullbacks,
fixed-order beta-function comparison, Green--Schwarz structure, heterotic
compactification targets, and circle momentum--winding exchange remain useful
parts of the proposed bridge when their missing quantum and global data are
supplied.

\item[Remaining boundary]
MTT has not yet selected a complete perturbative worldsheet theory on the
physical \(q=79\) visible--hidden background.  In the current research ledger,
blocker B.QG.01 remains open and depends on the physical Hull--Strominger
endpoint and the quantum BV/renormalized-transport layer.
\end{description}

\tableofcontents

\section{The corrected question}
\label{sec:intro}

\subsection{Target fields are necessary but not sufficient}

A conventional nonlinear sigma model starts with a map
\[
X:\Sigma\longrightarrow M
\]
and target couplings such as a metric \(G\), a two-form \(B\), and a dilaton
\(\Phi\).  In a heterotic model one also needs left-moving gauge degrees of
freedom and a gauge connection \(A\).  Pullback along \(X\) then produces the
familiar local couplings.  This is an important bridge from target geometry
to a worldsheet Lagrangian, but it does not determine the complete quantum
theory.

The field content, chirality, current algebra, normalization, gauge fixing,
ghost systems, physical state conditions, spin-structure sum, and
renormalization prescription are additional choices.  At higher genus one
also needs a measure on moduli space, modular covariance, sewing and
factorization, and treatment of degenerations.  Global anomaly and gerbe data
cannot be reconstructed from the local Lagrangian alone
\cite{Polchinski1998,GSW1987,Witten2012}.

The corrected MTT question is therefore:
\begin{quote}
Does one selected upper MTT state emit a complete, mutually compatible
perturbative-string record, and does the MTT evaluation of observables factor
through the standard worldsheet construction?
\end{quote}
This question can be answered precisely.  It is stronger than matching a few
target fields and weaker than claiming a nonperturbative definition of string
theory.

\subsection{Three claims that must not be merged}

We distinguish the following levels.
\begin{description}[leftmargin=3.0cm,style=nextline]
\item[Local encoding]
A map sends MTT variables or perturbations to worldsheet couplings.  The
companion proto-spinor paper owns a conditional quadratic version of this
statement \cite{NeroProtoWorldsheet2026}.

\item[Perturbative realization]
One complete, anomaly-free, modular and factorizing worldsheet record is
emitted, and the declared MTT observables factor through its amplitude map.
This is the conditional theorem of the present paper.

\item[Physical selection]
An MTT source, preparation, flow, or branch law chooses that record rather
than another allowed record.  Existence and consistency do not imply this
selection step \cite{NeroHeteroticSelection2026}.
\end{description}

\subsection{Argument map}

\Cref{sec:record} defines the complete lower object.
\Cref{sec:map} states the typed MTT encoding contract.
\Cref{sec:beta,sec:quantum-gates} separate target equations, central charge,
BRST, anomaly, modularity, and duality.
\Cref{sec:q79} audits the current selected \(q=79\) branch.
\Cref{sec:uv,sec:boundary} state exactly what fixed-genus ultraviolet
inheritance would require and what remains open.

\section{The lower object to be realized}
\label{sec:record}

\subsection{Target and local worldsheet data}

For definiteness, the main application is a closed oriented heterotic
worldsheet.  Other string theories require altered matter and projection
rows, but the logical separation is the same.

\begin{definition}[Perturbative-string record]
\label{def:string-record}
A perturbative-string record at a declared order \(N\) in \(\alpha'\), genus
range \(\cG\), and external-state class \(\cE\) is a tuple
\[
\mathfrak S=
(\mathfrak T,\mathfrak W,\mathfrak Q,\mathfrak P,
 \mathfrak A,\mathfrak M,\mathfrak I),
\]
with the following typed rows.
\begin{enumerate}[leftmargin=1.7em]
\item \(\mathfrak T\) is target data:
      \((M,G,B,\Phi,A,\nabla,\cG_B)\), including a global \(B\)-field or
      gerbe \(\cG_B\), a declared tangent connection \(\nabla\), and all
      compactification and bundle data used by the model.
\item \(\mathfrak W\) specifies the two-dimensional fields, chiralities,
      current algebra, local action, couplings, boundary conditions, and
      normalization.
\item \(\mathfrak Q\) specifies gauge fixing, matter and ghost Hilbert or
      state spaces, the BRST operator, its domain, and the physical state
      cohomology.
\item \(\mathfrak P\) specifies spin structures, GSO or analogous
      projections, and any orbifold or discrete-torsion phases.
\item \(\mathfrak A\) certifies local and global gauge, gravitational, and
      worldsheet anomalies, the differential Bianchi identity, and the
      perturbative order of the target equations.
\item \(\mathfrak M\) supplies genus-\(g\) integrands for
      \(g\in\cG\), their modular transformation law, the integration cycle or
      measure, and sewing and factorization on every relevant degeneration.
\item \(\mathfrak I\) supplies the tadpole, vacuum-shift, infrared, and soft
      prescription needed for the amplitudes in \(\cE\).
\end{enumerate}
\end{definition}

This definition does not demand an all-genus sum.  A record may be complete
for one fixed genus and multiplicity while an all-genus or nonperturbative
completion remains open.  The declared scope is part of the mathematical
object.

\subsection{The local sigma-model row}

For supplied target fields, the bosonic terms take the schematic form
\begin{align}
S_{\sigma}
={}&\frac{1}{4\pi\alpha'}\int_\Sigma
\sqrt{h}\,h^{ab}G_{MN}(X)\,
\partial_aX^M\partial_bX^N\,\dd^2\sigma \nonumber\\
&+\frac{\mathrm i}{4\pi\alpha'}\int_\Sigma
\epsilon^{ab}B_{MN}(X)\,
\partial_aX^M\partial_bX^N\,\dd^2\sigma
+\frac{1}{4\pi}\int_\Sigma
\sqrt{h}\,\Phi(X)R^{(2)}\,\dd^2\sigma ,
\label{eq:sigma-action}
\end{align}
after choosing Euclidean conventions.  A heterotic completion adds
right-moving worldsheet fermions, left-moving gauge degrees of freedom, and
their couplings.  Gauge fixing adds the appropriate ghost systems.

Pullback explains how a supplied target tensor enters
\cref{eq:sigma-action}.  It does not select the worldsheet field content,
prove quantum Weyl invariance, construct the GSO sum, or define the measure
on moduli space.  Consequently a bounded MTT projector can participate in an
encoding map, but boundedness alone cannot produce the RNS or heterotic
quantum theory.

\subsection{Compactification adds further rows}

For a heterotic compactification on a complex threefold \(X\), the target row
must include visible and hidden holomorphic bundles and connections in one
common stability chamber.  In a torsional \(SU(3)\)-structure background,
the Hull--Strominger equations include
\[
\dd(e^{-2\Phi}\omega^2)=0,\qquad
F^{0,2}=0,\qquad F\wedge\omega^2=0,\qquad
H=\mathrm i(\bar\partial-\partial)\omega,
\]
together with a differential Green--Schwarz identity.  Solving these
target-space equations at first order in \(\alpha'\) is not by itself an
exact \((0,2)\) SCFT \cite{Strominger1986,FuYau2008}.

The companion Hull--Strominger paper gives the appropriate conditional
fixed-point correspondence; the companion compactification audit records
which older bundle examples fail \cite{NeroMTTStrominger2026,NeroFluxAudit2026}.
Those theorems are imported here rather than duplicated.

\section{The MTT encoding contract}
\label{sec:map}

\subsection{Typed source and evaluation maps}

Let \(\cU_{\mathrm{MTT}}\) be a declared upper MTT configuration space.  A
string encoding is a partial typed map
\[
\cR_{\mathrm{str}}:
\cU_{\mathrm{MTT}}\dashrightarrow\cM_{\mathrm{str}},
\]
where \(\cM_{\mathrm{str}}\) is the space of records in
\cref{def:string-record}.  ``Typed'' means that each lower field, bundle,
connection, projection, character, and measure has an identified upper
source and convention.  Matching dimensions or abstractly isomorphic groups
do not define this map.

For a selected scope \((N,\cG,\cE)\), let
\[
\cA_{\mathrm{ws}}:\cM_{\mathrm{str}}\dashrightarrow\cO_{\cG,\cE}
\]
be the standard perturbative worldsheet evaluation map, and let
\[
\cA_{\mathrm{MTT}}:\cU_{\mathrm{MTT}}\dashrightarrow\cO_{\cG,\cE}
\]
be the corresponding MTT evaluation.  The required commuting relation is
\[
\cA_{\mathrm{MTT}}
=\cA_{\mathrm{ws}}\circ\cR_{\mathrm{str}}
\quad\text{on a declared domain }\cU_0.
\label{eq:factorization}
\]

\begin{definition}[Complete MTT string realization]
\label{def:mtt-realization}
An upper state \(u_\ast\in\cU_0\) is a complete MTT realization of a
perturbative string sector at scope \((N,\cG,\cE)\) if:
\begin{enumerate}[leftmargin=1.7em]
\item \(\cR_{\mathrm{str}}(u_\ast)\) supplies every row of
      \cref{def:string-record};
\item those rows satisfy their standard compatibility, anomaly, BRST,
      modular, factorization, and infrared conditions at the declared scope;
\item every approximation and field-redefinition convention is explicit;
      and
\item \cref{eq:factorization} holds at \(u_\ast\).
\end{enumerate}
\end{definition}

\begin{theorem}[Conditional perturbative-string realization]
\label{thm:string-realization}
If \(u_\ast\) satisfies \cref{def:mtt-realization}, then
\(\cR_{\mathrm{str}}(u_\ast)\) is a consistent perturbative-string record at
scope \((N,\cG,\cE)\), and every observable in that scope obeys
\[
\cA_{\mathrm{MTT}}(u_\ast)
=\cA_{\mathrm{ws}}\!\left(\cR_{\mathrm{str}}(u_\ast)\right).
\]
\end{theorem}

\begin{proof}
The first three hypotheses put the emitted record in the domain of the
standard amplitude map with a fixed approximation and infrared convention.
The last hypothesis is exactly the displayed equality.  The conclusion does
not add an existence, uniqueness, selection, or all-genus assertion beyond
the supplied record.
\end{proof}

\begin{remark}[Why a conditional theorem is useful]
The theorem is not intended to hide the hard work in an assumption.  It
locates that work.  A proposed MTT construction can now be tested row by row,
and no target-space computation can silently stand in for a missing GSO,
modular, BV, or infrared certificate.
\end{remark}

\subsection{Relation to the proto-spinor bridge}

The proto-spinor/worldsheet paper studies a different arrow.  Its conditional
quadratic bridge compares Hessians near one background:
\[
L^\ast H_{\mathrm{ws}}L=H_{\mathrm{ps}}
\]
on a selected tangent subspace, with a cubic remainder estimate
\cite{NeroProtoWorldsheet2026}.  That theorem can provide a local component
of \(\cR_{\mathrm{str}}\), but it does not emit the global record in
\cref{def:string-record}.  Conversely, the present theorem does not re-prove
the local Hessian statement.  The two papers have complementary ownership.

\section{Beta functions and target equations}
\label{sec:beta}

\subsection{What the standard comparison says}

Renormalization of a nonlinear sigma model produces beta functions on the
space of target couplings.  To leading order, in a common convention,
\begin{align}
\beta^G_{MN}
&=\alpha'\left(
R_{MN}-\frac14H_{MPQ}H_N{}^{PQ}
+2\nabla_M\nabla_N\Phi
\right)+O(\alpha'^2), \\
\beta^B_{MN}
&=\alpha'\left(
-\frac12\nabla^PH_{PMN}
+\nabla^P\Phi\,H_{PMN}
\right)+O(\alpha'^2),
\end{align}
with gauge and dilaton terms determined by the chosen string model and
scheme.  Vanishing Weyl-anomaly coefficients yields target field equations
order by order, up to local field redefinitions
\cite{Friedan1980,CallanEtAl1985,FradkinTseytlin1985}.

This comparison has two important qualifications.
\begin{enumerate}[leftmargin=1.7em]
\item A solution through order \(\alpha'^N\) is not an all-orders fixed
      point and does not by itself construct an exact CFT.
\item Beta functions, Weyl-anomaly coefficients, and target effective-action
      equations are related in a specified renormalization scheme; they are
      not literally the same coordinate-free object.
\end{enumerate}

\subsection{Controlled residual transfer}

\begin{proposition}[Fixed-order residual transfer]
\label{prop:residual-transfer}
Let \(E_N(u)\) be the vector of target effective-action residuals emitted by
an MTT state \(u\), and let \(\beta_N(u)\) be the corresponding vector of
worldsheet Weyl-anomaly coefficients through order \(\alpha'^N\).  Suppose
on a neighborhood \(U\) that
\[
\beta_N(u)=K_N(u)E_N(u)+r_{N+1}(u),
\]
where \(K_N(u)\) is bounded by \(C_K\) in the chosen norms and
\(\|r_{N+1}(u)\|\le C_r|\alpha'|^{N+1}\).  Then
\[
\|\beta_N(u)\|
\le C_K\|E_N(u)\|+C_r|\alpha'|^{N+1}
\qquad (u\in U).
\]
In particular, an exact projected target solution through order \(N\)
implies only an \(O(\alpha'^{N+1})\) Weyl-anomaly residual under these
hypotheses.
\end{proposition}

\begin{proof}
Take norms in the assumed relation and apply the operator-norm bound and the
triangle inequality.
\end{proof}

The proposition is deliberately an error statement.  To use it, MTT must
emit the same target fields and scheme used by \(K_N\), and the residual
bound must include every active equation.  It cannot promote a first-order
Hull--Strominger solution to an exact infrared SCFT.

\section{Independent quantum consistency gates}
\label{sec:quantum-gates}

\subsection{Criticality and ghosts}

The central-charge check is exact for the standard free heterotic field
content, but it is a consistency calculation rather than a derivation of that
content.

\begin{lemma}[Standard heterotic central-charge arithmetic]
\label{lem:central-charge}
For ten target coordinate bosons, sixteen units of left-moving gauge current
central charge, ten right-moving Majorana fermions, and the standard
reparametrization and superconformal ghosts,
\[
c_L^{\mathrm{tot}}=(10+16)-26=0,\qquad
c_R^{\mathrm{tot}}=\left(10+\frac{10}{2}\right)-15=0.
\]
\end{lemma}

\begin{proof}
A free real boson contributes \(1\), a free real Majorana fermion contributes
\(1/2\), the left-moving \(bc\) ghosts contribute \(-26\), and the combined
right-moving \(bc\) and \(\beta\gamma\) ghosts contribute \(-15\).  Summing
the stated standard field content gives the result
\cite{GSW1987,Polchinski1998}.
\end{proof}

This calculation checks the critical bookkeeping once the heterotic matter
and ghost systems have been chosen.  It does not show that an MTT projection
selects ten coordinates, the gauge current algebra, or the GSO projection.

\subsection{BRST and physical states}

BRST nilpotence requires cancellation of the relevant quantum anomalies and
a correctly defined operator on the gauge-fixed state space.  Target
beta-function equations are part of the consistency conditions for a
background, but the slogan
\[
\text{``BRST nilpotence''}\quad\Longleftrightarrow\quad\text{``leading
target beta functions vanish''}
\]
is too coarse.  One must specify the matter-plus-ghost theory, normal
ordering, central charge, background, and perturbative order.  The physical
states are then identified by BRST cohomology, subject to the chosen
projection and inner product.

\subsection{Spin structures and GSO}

The GSO projection is not a sign appended after the target equations.  It
controls the physical spectrum, spacetime spin-statistics, tachyon removal,
and modular consistency.  At genus one and beyond, the spin-structure sum
must transform correctly and be compatible with factorization.  The original
GSO construction supplies the standard supersymmetric example
\cite{GSO1976}; a selected \(q=79\) model must still emit its own analytic
characters and phases.

\subsection{Green--Schwarz and global B-field data}

In a heterotic convention one writes
\[
H=\dd B-\frac{\alpha'}4
\bigl(\omega_3(A)-\omega_3(\nabla)\bigr),
\qquad
\dd H=\frac{\alpha'}4
\left(\tr R_\nabla\wedge R_\nabla-\tr F\wedge F\right),
\label{eq:bianchi}
\]
with any visible, hidden, and five-brane contributions made explicit.
Green--Schwarz factorization is a ten-dimensional anomaly statement
\cite{GreenSchwarz1984}.  A compactification additionally needs a
differential representative satisfying \cref{eq:bianchi} and compatible
global gerbe data.

Equality of integrals over all four-cycles establishes equality of the
relevant de Rham cohomology classes under appropriate hypotheses.  It does
not make two chosen curvature four-forms equal pointwise, construct \(H\), or
trivialize a torsion differential-cohomology obstruction.  Freed--Witten
conditions and related global restrictions are likewise separate from local
curvature arithmetic \cite{FreedWitten1999}.

\subsection{Modularity and factorization}

For a modular-invariant closed-string integrand, integration may be
restricted to a fundamental domain of the torus modular group.  The region
that resembles arbitrarily short proper time in a point-particle
representation is then identified with another region; the remaining
boundary at large imaginary modulus is associated with degeneration and
factorization, where infrared and tadpole issues can occur
\cite{Witten2012,SenUVIR2016}.

This mechanism cannot be invoked before modular invariance, the measure,
spin-structure sum, and factorization have been established for the actual
model.  A finite set of modularly permuted labels is useful algebraic input,
but it is not yet an analytic partition function.

\section{Duality: what the circle calculation proves}
\label{sec:duality}

For a compact free boson on a circle, the zero-mode mass contribution is
schematically
\[
M^2_{\mathrm{zero}}
=\frac{n^2}{R^2}+\frac{w^2R^2}{\alpha'^2}.
\]
It is invariant under
\[
(n,w,R)\longmapsto(w,n,\alpha'/R).
\]
This exact arithmetic is a valuable compatibility witness for momentum and
winding exchange.

Full quantum T-duality contains more.  In a nonlinear sigma model it requires
an isometry or another specified duality structure, transformation of
\((G,B,\Phi)\), the functional measure and dilaton shift, global and
topological data, and compatibility with the projection and spectrum
\cite{Buscher1988}.  Therefore an MTT shared-circle spectrum can encode the
kinematic seed of T-duality without yet proving equality of two complete
string backgrounds.

S-duality, holography, and an M-theory lift require additional nonperturbative
or boundary data and are outside the theorem proved here.  Their possible
MTT encodings belong in separate papers with their own source maps; they are
not consequences of \cref{thm:string-realization}.

\section{The selected q=79 worldsheet cutset}
\label{sec:q79}

\subsection{Why this is the relevant branch}

The current selected MTT route toward a heterotic string realization uses
the time-oriented \(q=79/F\) branch and a Fu--Yau-type non-K\"ahler
compactification target.  It is not the strict Calabi--Yau branch described
in the companion paper \cite{NeroCalabiYau2026}.  The exact finite arithmetic,
Cech carrier, and projected HYM calculations give substantial support at
their declared tiers, but only the worldsheet cutset decides whether they
assemble into the record of \cref{def:string-record}.

\begin{table}[htbp]
\centering
\small
\begin{tabularx}{\textwidth}{@{}p{0.05\textwidth}p{0.22\textwidth}p{0.13\textwidth}X@{}}
\toprule
Row & Required object & State & Current evidence or missing endpoint \\
\midrule
W1 & Selected target branch & Available &
Time-oriented \(q=79/F\) representative. \\
W2 & Charge and Bianchi sector & Available &
Fu--Yau/Mukai charge data and Green--Schwarz Bianchi sector at the declared
topological and curvature tiers. \\
W3 & Visible local anomaly row & Available &
Curvature-level visible Green--Schwarz cancellation.  This is not the full
global gerbe or physical visible--hidden bundle. \\
W4 & Critical central charge & Available &
The universal heterotic arithmetic in \cref{lem:central-charge}.  It checks
standard field content rather than selecting it. \\
W5 & Low-energy limit & Available &
The \(q=79\) GR and quantum-EFT limit at the declared parity/controlled tier.
It is not a worldsheet completion. \\
\bottomrule
\end{tabularx}
\caption{Available rows in the current \(q=79\) heterotic worldsheet
contract.  Availability is limited to the stated row and tier.}
\label{tab:q79-available}
\end{table}

\begin{table}[htbp]
\centering
\small
\begin{tabularx}{\textwidth}{@{}p{0.05\textwidth}p{0.22\textwidth}p{0.13\textwidth}X@{}}
\toprule
Row & Required object & State & Current evidence or missing endpoint \\
\midrule
W6 & Global differential data & Open &
Full Deligne gerbe, Freed--Witten restrictions, and differential class on the
complete visible cycle set. \\
W7 & All-order target background & Open &
The current flux construction is first order in \(\alpha'\); an all-order
background or a controlled exact alternative is absent. \\
W8 & Exact heterotic \((0,2)\) SCFT & Partial &
K3 GLSM, local anomaly, topological clutching candidates, and related finite
data exist.  The physical non-pullback holomorphic visible--hidden bundle,
HYM connection, differential Bianchi representative, GSO currents, and exact
IR SCFT remain open. \\
W9 & Modular/GSO amplitudes & Partial &
The finite \(F_3^2\) torsion phase, seven modular seed orbits, and finite
covariance are exact.  Oscillator and gauge-current characters,
spin-structure sum, multipliers, analytic GSO completion, and factorization
are missing. \\
W10 & Quantum string-field/BV data & Open &
No \(q=79\)-specific vertices and quantum BV master-action certificate have
been emitted. \\
W11 & Tadpole and infrared control & Open &
Vacuum shift, tadpole cancellation, massless soft behavior, and degeneration
prescriptions remain to be supplied on the same branch. \\
W12 & Beyond fixed genus & Open &
No all-genus summability or nonperturbative definition at finite string
coupling is claimed. \\
\bottomrule
\end{tabularx}
\caption{Partial and open rows in the current \(q=79\) heterotic worldsheet
contract.}
\label{tab:q79-open}
\end{table}

The count is therefore
\[
\boxed{5\ \text{available},\qquad 2\ \text{partial},\qquad 5\ \text{open}.}
\]
Rows W8 and W9 are genuine advances: they reduce the unknown worldsheet data
to a much sharper finite and analytic interface.  They do not count as
complete until their missing objects are constructed on the same physical
background.

\subsection{The physical-bundle dependency}

The principal geometric dependency is not another four-dimensional fit.  It
is one physical non-pullback visible--hidden Hull--Strominger endpoint with
holomorphic bundles, a common positive HYM chamber, connections, and the
differential anomaly identity.  The current topological clutching data show
that the required instanton and chirality targets are not excluded, but
topological admissibility does not prove holomorphic existence or HYM.

This is why the \(q=79\) worldsheet blocker depends on the physical
Hull--Strominger endpoint.  The quantum amplitude side independently depends
on a renormalized BV/QME transport or an equivalent construction.  Both
dependencies must be discharged before the present conditional theorem can
be instantiated.

\subsection{Finite modular data: exact but incomplete}

The selected finite gerbe cocycle gives an exact discrete-torsion phase on 81
torus twist sectors.  Their modular \(S,T\) action reduces to seven seed
blocks, and the associated twisted group algebra is
\(\operatorname{Mat}_3(\mathbb C)\).  These are exact finite statements.

Analytic characters carry oscillator spectra, current algebras, spin
structures, and multiplier systems.  The finite permutation/cocycle layer
constrains how such characters may transform, but it cannot supply their
\(\tau\)-dependence or prove factorization.  This distinction prevents a
finite modular covariance certificate from being mistaken for a one-loop
string partition function.

\section{Ultraviolet inheritance at fixed genus}
\label{sec:uv}

The current q79 research corpus contains the following conditional result,
which is imported rather than re-owned here:
\begin{quote}
If one same-source \(q=79/F\) background supplies an exact anomaly-free
modular heterotic \((0,2)\) SCFT, a tachyon-free GSO projection,
factorization, a heterotic quantum BV master action, and a tadpole/infrared
prescription, then fixed-genus, fixed-multiplicity amplitudes have no local
ultraviolet divergences.
\end{quote}
The logic is standard string perturbation theory: local short-tube regions
are controlled by modular equivalence, while boundaries of moduli space are
handled through factorization and infrared prescriptions
\cite{Witten2012,SenBV2015,SenUVIR2016}.

This result is not a permanent Gaussian damping of the physical graviton
propagator.  It is conditional inheritance from a complete worldsheet.
Because W8--W11 are not closed, it is not yet an unconditional \(q=79\)
ultraviolet-completion theorem.  Even after those fixed-genus rows close,
convergence of the sum over genera and a nonperturbative definition remain
the independent W12 boundary.

\section{Four-dimensional reduction and prediction}
\label{sec:eft}

Once a complete compactification record is supplied, standard dimensional
reduction can produce a four-dimensional effective action.  Gauge couplings,
Yukawa couplings, masses, and threshold corrections depend on normalized
internal modes, moduli, bundle data, and a renormalization and matching
scheme.  For example, a holomorphic Yukawa is schematically an internal
overlap integral, but its physical value also depends on kinetic
normalization and transport to the comparison scale.

The finite \(27\times27\) MTT carrier and the current Standard Model parity
and profile calculations can be compatible with such a reduction.  They do
not become string-derived merely because the same finite ranks appear.
Conversely, a future worldsheet construction would not turn measured profile
inputs into source-selected predictions unless its normalized overlap
kernels emit those values independently.

The factorization required for a genuine string origin is therefore a
commuting chain
\[
\begin{split}
u_\ast
&\xmapsto{\ \cR_{\mathrm{str}}\ }
\mathfrak S_{q79}
\xmapsto{\ \text{compactification}\ }
\mathfrak E_4\\
&\xmapsto{\ \text{RG and matching}\ }
\cO_{\mathrm{physical}},
\end{split}
\]
with the same source, conventions, and uncertainty budget throughout.  Each
arrow is an independent theorem or certified computation.

\section{Claim audit and completion plan}
\label{sec:boundary}

\subsection{Disposition of Version 1 claims}

\begin{center}
\small
\begin{tabularx}{0.96\textwidth}{@{}p{0.40\textwidth}p{0.16\textwidth}X@{}}
\toprule
Former claim & Status & Version 2 resolution \\
\midrule
Bounded projection produces Polyakov/RNS theory & Withdrawn &
Pullback produces local couplings only after the worldsheet fields and action
are supplied; quantum data remain separate. \\
An MTT fixed point implies \(\beta=0\) & Conditional &
Replaced by \cref{prop:residual-transfer} with scheme, order, completeness,
and error hypotheses. \\
BRST nilpotence is equivalent to leading beta functions & Withdrawn &
BRST requires the full matter-plus-ghost quantum system and anomaly
cancellation. \\
Characteristic-number equality proves the Bianchi identity & Withdrawn &
Cohomological equality is not a differential or global gerbe certificate. \\
The Hull--Strominger solution is the unique MTT minimizer & Withdrawn &
The adjacent paper now proves only a conditional fixed-point correspondence. \\
Modular invariance is automatic in the compactification corner & Withdrawn &
The actual characters, GSO sum, measure, sewing, and factorization must be
constructed. \\
Circle spectral symmetry proves exact T-duality & Narrowed &
It proves the momentum--winding seed; full Buscher and quantum data are
additional. \\
MTT has a complete ultraviolet-finite string sector & Conditional &
Fixed-genus UV inheritance follows only after the W8--W11 worldsheet
contract is supplied. \\
Curvature-gap dynamics fix the four-dimensional potential & Withdrawn &
Moduli stabilization and normalized reduction require an independent source
and calculation. \\
S-duality, holography, and M-theory follow from this bridge & Out of scope &
They require separate nonperturbative or boundary realization maps. \\
\bottomrule
\end{tabularx}
\end{center}

\subsection{The shortest route to an instantiated theorem}

The present theorem becomes a physical \(q=79\) result only after the
following same-source sequence is completed.
\begin{enumerate}[leftmargin=1.7em]
\item Construct the non-pullback visible and hidden holomorphic bundles on
      the selected Fu--Yau carrier, their common HYM chamber and connections,
      and the differential/global anomaly data.
\item Derive or construct the exact infrared \((0,2)\) SCFT and identify the
      target fields with the physical bundle rather than an aggregate local
      anomaly model.
\item Emit the seven analytic seed characters, oscillator and current
      sectors, spin-structure and GSO sum, multipliers, modular mixing, and
      factorization.
\item Construct the \(q=79\)-specific string-field vertices or an equivalent
      quantum BV/QME certificate.
\item Supply the tadpole, vacuum-shift, soft, and infrared prescription.
\item Evaluate the factorization
      \(\cA_{\mathrm{MTT}}=\cA_{\mathrm{ws}}\circ\cR_{\mathrm{str}}\) on a
      nontrivial held-out observable set.
\end{enumerate}
All-genus summability or a nonperturbative definition is a subsequent
question, not a hidden seventh item in the fixed-genus claim.

\subsection{Conclusion}

String theory remains a serious and promising realization language for MTT,
especially on the selected \(q=79\) Fu--Yau branch.  The corrected relation is
now exact in its logic.  MTT may supply an upper source and selection
mechanism; standard worldsheet theory supplies the consistency and amplitude
machinery; a typed source map and commuting evaluation diagram must connect
them.

Five of the twelve \(q=79\) worldsheet rows are available and two are
partially constructed.  This is meaningful progress, but it is not a
complete string background.  The remaining work is no longer summarized by
the vague phrase ``derive the worldsheet.''  It is the concrete physical
bundle, IR SCFT, analytic modular/GSO, BV, and infrared packet listed above.
When those rows are emitted from one source, \cref{thm:string-realization}
will turn the present conditional encoding into an instantiated
perturbative-string result.

\bibliographystyle{plain}

% BEGIN MTT MANAGED COMPUTATIONAL EVIDENCE
\section*{Computational Evidence and Reproducibility}
The exact q79 arithmetic, finite rank-two Cech witness, and rank-two Wiener-contraction certificate are contextual evidence for the selected heterotic branch. They do not directly construct the physical non-pullback visible-hidden bundle, exact infrared (0,2) SCFT, analytic modular/GSO characters, quantum BV packet, tadpole prescription, all-genus completion, or the typed MTT-to-worldsheet evaluation factorization.

The referenced rows are frozen to the curated results repository state identified below. Hashes are grouped in eight-character blocks for line breaking.
\begin{quote}\small
\textbf{Repository:} \url{https://github.com/PeterNero/mtt-results-repro}\\
\textbf{Commit:} \texttt{31247ebb 5c22f3fb b5443024 365433c6 ee0bff4a}\\
\textbf{Manifest:} \path{release/result_manifest.json}\\
\textbf{Manifest SHA-256:}\\
\texttt{fb399689 60b00584 631dbf53 1a708e18}\\
\texttt{ef928d6b 6d935119 c185d7f6 32b1e7cd}
\end{quote}

Tier labels are quoted verbatim from that manifest. A row used directly supports only the specific computational statement identified above; a corpus-state cross-check does not prove this paper's local theorems; and an open row is evidence of an unresolved obligation, never of closure.

\par\begingroup\dimen0=\pagegoal\advance\dimen0 by -\pagetotal\ifdim\dimen0<7\baselineskip\newpage\fi\endgroup
\paragraph{Corpus-state cross-checks.}
\begin{itemize}
\setlength{\itemsep}{0.25em}
\item \path{A19/hym_wiener_contraction} (\emph{numeric certified}).\par Weighted-theta Fourier-tail and Wiener contraction certificate.
\item \path{A07/literal_cech_witness} (\emph{derived exact}).\par Literal 81-entry, 729-cocycle finite Cech witness.
\item \path{A11/q79_exact_audit} (\emph{derived exact}).\par Executable q=79 exact-branch audit.
\item \path{A11/q79_exact_theorem} (\emph{derived exact}).\par CRT q=79 theorem on the selected exact branch.
\end{itemize}

No imported row changes theorem ownership or promotes a neighboring claim: all local statements retain their stated hypotheses, domains, and limitations.
% END MTT MANAGED COMPUTATIONAL EVIDENCE

\bibliography{main}

\end{document}
